% !TeX root = ../../../main.tex
\subsubsection{Shared Forecasting Sensitivity Analyses}\label{ch07:sec:shared-sensitivity-results}
\ChSevenHeadingOnly

\paragraph{Time-Ordered versus Text-Ordered Training}\label{ch07:apd:sensitivity_analyses}

{\raggedright
Table \ref{ch07:tab:dsr1q_l33_comparison} shows time-ordering generally improves concordance, particularly for \texttt{DeepSeek-R1} annotations, while F1 scores show mixed results. Text-ordering sometimes yields better external dataset performance with \texttt{Llama-3.3-70B}, suggesting that preserving the narrative sequence can support accurate classification even when it does not follow temporal order.  These results highlight a trade-off between temporal alignment and narrative coherence for different tasks.


\par}
% !TeX root = ../../../main.tex
\begin{table}[htbp]
\caption{Comparison of text-ordered and time-ordered training for \texttt{DeepSeek-R1} and \texttt{Llama 3.3-70B} annotations on F1 scores and c-index values. 
}\ChSevenFloatLabel{ch07:tab:dsr1q_l33_comparison}
\ChSevenTableBodyStart
\begin{ChSevenFloatBody}

\centering
{
\begin{adjustbox}{max width=\linewidth}
\begin{tabular}{lcccccccc}
\toprule
& F1(1h) & F1(1d) & F1(1w) & c-index & F1(1d)-10 & c10 & F1(1d)-100 & c100 \\
\hline
\multicolumn{9}{c}{\textbf{\textit{T2S2} (\texttt{DeepSeek-R1})}} \\
\hline
\textbf{Text-ordered} & 0.290 & 0.569 & 0.876 & 0.659 & 0.361 & 0.598 & 0.615 & 0.609 \\
\textbf{Time-ordered} & 0.306 & 0.576 & 0.877 & 0.672 & 0.428 & 0.576 & 0.607 & 0.613 \\
\hline
\multicolumn{9}{c}{\textbf{\textit{T2S2} (\texttt{Llama-3.3-70B})}} \\
\hline
\textbf{Text-ordered} & 0.311 & 0.648 & 0.898 & 0.586 & 0.423 & 0.526 & 0.662 & 0.541 \\
\textbf{Time-ordered} & 0.257 & 0.645 & 0.903 & 0.632 & 0.382 & 0.595 & 0.616 & 0.612 \\
\bottomrule
\end{tabular}
\end{adjustbox}
}



\end{ChSevenFloatBody}
\end{table}

\paragraph{Masking History}\label{ch07:apd:masking_l33_dsr1}

Tables~\ref{ch07:tab:masking_effects_DS} and~\ref{ch07:tab:masking_effects_llama} report the effect of timestep dropout on classification and ranking performance.
We evaluate performance across three datasets (\textit{T2S2}, \textit{sepsis-10}, and \textit{sepsis-100}) using F1 scores for prediction of next \(k\) events happening within the next day and concordance index (c-index) for correctly ordering the next \(k\) events.
As the proportion of masked history increases, we observe a general decline in F1 performance, particularly at higher masking rates (e.g., 90\%), suggesting that access to historical context is important for accurate forecasting.
Interestingly, the concordance index remains relatively stable, indicating the model's robustness in ranking the next \(k\) events with partial input histories.

More specifically, we can see that F1 scores generally decrease with increasing TDR, with relatively mild degradation up to 30--60\% dropout and sharper declines thereafter.
For example, in the \texttt{DeepSeek-R1} annotations (Table~\ref{ch07:tab:masking_effects_DS}), \textit{T2S2}/\textit{sepsis-10}/\textit{sepsis-100} drop from 0.576/0.428/0.607 at TDR = 0\% to 0.423/0.200/0.510 at TDR = 90\%, with \textit{sepsis-10} showing the largest relative decline.
In contrast, concordance (c-index) remains stable across dropout levels (0.661--0.678), indicating that ranking-based evaluations are less sensitive to history truncation.
These trends are consistent for \texttt{Llama-3.3} annotations (Table~\ref{ch07:tab:masking_effects_llama}), suggesting that ranking objectives capture temporal relationships that are more robust to missing historical information than binary classification accuracy of predicting whether an event happens within a time window.


% !TeX root = ../../../main.tex
\begin{table}[htbp]
\caption{Effect of Randomly Masking History (\texttt{DeepSeek-R1} annotations) on F1 and Concordance. \textbf{Abbreviations}: F1(1d)-10/100: F1 (1 day) for \textit{sepsis-10} and \textit{sepsis-100} respectively. c10/100: c-index for \textit{sepsis-10} and \textit{sepsis-100} respectively.}\ChSevenFloatLabel{ch07:tab:masking_effects_DS}
\ChSevenTableBodyStart
\begin{ChSevenFloatBody}

    \centering
    \begin{tabular}{lcccccc}
        \toprule
        \textbf{Timestep drop rate} & F1(1d) & c-index & F1(1d)-10 & c10 & F1(1d)-100 & c100 \\
        \hline
        0\%  & 0.576 & 0.672 & 0.428 & 0.576 & 0.607 & 0.613 \\
        15\% & 0.572 & 0.667 & 0.415 & 0.564 & 0.604 & 0.616 \\
        30\% & 0.548 & 0.672 & 0.357 & 0.542 & 0.591 & 0.619 \\
        60\% & 0.542 & 0.678 & 0.310 & 0.578 & 0.588 & 0.616 \\
        90\% & 0.423 & 0.661 & 0.200 & 0.561 & 0.510 & 0.614 \\
        \bottomrule
    \end{tabular}
    
    

\end{ChSevenFloatBody}
\end{table}



% !TeX root = ../../../main.tex
\begin{table}[htbp]
\caption{Effect of Randomly Masking History (\texttt{Llama-3.3-70B} annotations) on F1 and Concordance. \textbf{Abbreviations}: F1(1d)-10/100: F1 (1 day) for \textit{sepsis-10} and \textit{sepsis-100} respectively. c10/100: c-index for \textit{sepsis-10} and \textit{sepsis-100} respectively.}\ChSevenFloatLabel{ch07:tab:masking_effects_llama}
\ChSevenTableBodyStart
\begin{ChSevenFloatBody}

    \centering
    \resizebox{0.75\linewidth}{!}{
    \begin{tabular}{lcccccc}
        \toprule
        \textbf{Timestep drop rate} & F1(1d) & c-index & F1(1d)-10 & c10 & F1(1d)-100 & c100 \\
        \hline
        0\%  & 0.645 & 0.632 & 0.431 & 0.595 & 0.626 & 0.612 \\
        15\% & 0.628 & 0.635 & 0.422 & 0.592 & 0.624 & 0.610 \\
        30\% & 0.634 & 0.634 & 0.423 & 0.586 & 0.621 & 0.598 \\
        60\% & 0.604 & 0.634 & 0.391 & 0.577 & 0.605 & 0.613 \\
        90\% & 0.531 & 0.622 & 0.328 & 0.546 & 0.534 & 0.607 \\
        \bottomrule
    \end{tabular}
    }
    
    

\end{ChSevenFloatBody}
\end{table}

\FloatBarrier
